% Part: methods
% Chapter: induction
% Section: induction-on-N

\documentclass[../../../include/open-logic-section]{subfiles}

\begin{document}

\olfileid[hi]{mth}{ind}{inN}

\olsection{$\Nat$ पर आगमन}

अपने सरलतम रूप में आगमन वह तकनीक है जिससे सभी प्राकृतिक संख्याओं के लिए
परिणाम सिद्ध किए जाते हैं। इसमें इस तथ्य का उपयोग होता है कि $0$ से आरंभ
करके बार-बार $1$ जोड़ने पर अंततः हम प्रत्येक प्राकृतिक संख्या तक पहुँचते
हैं। अतः यह सिद्ध करने के लिए कि कोई बात प्रत्येक संख्या के लिए सत्य है, हम
(1) स्थापित कर सकते हैं कि वह $0$ के लिए सत्य है और (2) दिखा सकते हैं कि
जब भी वह किसी संख्या $n$ के लिए सत्य हो, तो अगली संख्या $n+1$ के लिए भी
सत्य है। यदि ``संख्या $n$ में गुण $P$ है'' को $P(n)$ से (और ``संख्या $k$
में गुण $P$ है'' को $P(k)$ से, इत्यादि) संक्षिप्त करें, तो सभी
$n\in\Nat$ के लिए $P(n)$ का आगमन द्वारा प्रमाण इनसे बनता है:
\begin{enumerate}
\item $P(0)$ का प्रमाण, और
\item यह प्रमाण कि किसी भी $k$ के लिए यदि $P(k)$, तो $P(k+1)$।
\end{enumerate}
इसे बिल्कुल स्पष्ट करने के लिए मान लें कि हमारे पास (1) और (2) दोनों हैं।
तब (1) बताता है कि $P(0)$ सत्य है। यदि (2) भी हो, तो विशेषतः हम जानते हैं
कि यदि $P(0)$, तो $P(0+1)$, अर्थात् $P(1)$। यह सामान्य कथन ``किसी भी $k$
के लिए यदि $P(k)$, तो $P(k+1)$'' में $k$ के स्थान पर $0$ रखने से निष्पन्न
होता है। अतः मोडस पोनेन्स से हमारे पास $P(1)$ है। (2) से फिर, अब $n$ के
स्थान पर $1$ लेकर, हमारे पास है: यदि $P(1)$, तो $P(2)$। चूँकि हमने अभी
$P(1)$ स्थापित किया, मोडस पोनेन्स से हमारे पास $P(2)$ है। और इसी प्रकार
आगे। किसी भी संख्या $n$ के लिए ऐसा $n$ बार करने के बाद अंततः हम $P(n)$ पर
पहुँचते हैं। इसलिए (1) और (2) मिलकर किसी भी $n\in\Nat$ के लिए $P(n)$
स्थापित करते हैं।

एक उदाहरण देखें। मान लें हम जानना चाहते हैं कि $n$ पासों से कितने भिन्न
योग फेंके जा सकते हैं। कुछ मूर्खतापूर्ण लगने पर भी $0$ पासों से आरंभ करें।
यदि आपके पास कोई पासा नहीं है, तो केवल एक संभव योग ``फेंका'' जा सकता है:
कोई बिंदु नहीं, जिनका योग $0$ है। अतः भिन्न संभव फेंकों की संख्या $1$ है।
यदि केवल एक पासा हो, अर्थात् $n=1$, तो $1$ से $6$ तक छह संभव मान हैं। दो
पासों से $2$ से $12$ तक कोई भी योग फेंक सकते हैं, अर्थात् $11$ संभावनाएँ।
तीन पासों से $3$ से $18$ तक कोई भी संख्या फेंक सकते हैं, अर्थात् $16$
भिन्न संभावनाएँ। $1$, $6$, $11$, $16$: कोई प्रतिरूप दिखता है; शायद उत्तर
$5n+1$ है? निस्संदेह $5n+1$ अधिकतम संभव है, क्योंकि $n$ पासों से फेंके जा
सकने वाले न्यूनतम मान $n$ (सभी $1$) और अधिकतम मान $6n$ (सभी $6$) के बीच
केवल $5n+1$ संख्याएँ हैं।

\begin{thm}
  $n$ पासों से $n$ और $6n$ के बीच के सभी $5n+1$ संभव मान फेंके जा सकते
  हैं।
\end{thm}

\begin{proof}
$P(n)$ यह दावा हो: ``$n$ पासों से $n$ और $6n$ के बीच कोई भी संख्या फेंकना
संभव है।'' आगमन का उपयोग करने के लिए हम सिद्ध करते हैं:
\begin{enumerate}
\item \emph{आगमन आधार} $P(1)$, अर्थात् केवल एक पासे से $1$ और $6$ के बीच
  कोई भी संख्या फेंकी जा सकती है।
\item \emph{आगमन चरण}, सभी $k$ के लिए यदि $P(k)$, तो $P(k+1)$।
\end{enumerate}

(1) छह फलकों वाले पासे की जाँच से सिद्ध होता है। उसके सभी छह फलक हैं और
$1$ तथा $6$ के बीच प्रत्येक संख्या किसी एक फलक पर आती है। अतः एक पासे से
$1$ और $6$ के बीच कोई भी संख्या फेंकना संभव है।

(2) सिद्ध करने के लिए हम सशर्त के पूर्ववर्ती, अर्थात् $P(k)$, को मानते हैं।
इस मान्यता को \emph{आगमन परिकल्पना} कहते हैं। हम इसका उपयोग $P(k+1)$ सिद्ध
करने में करते हैं। कठिन भाग यह ढूँढना है कि $k+1$ पासों के फेंक के संभव
मानों को $k$ पासों के फेंकों के संभव मानों और अतिरिक्त $(k+1)$वें पासे के
फेंक के पदों में कैसे सोचा जाए---आगमन परिकल्पना का उपयोग करना हो तो हमें यही
करना होगा।

आगमन परिकल्पना कहती है कि $k$ पासों से $k$ और $6k$ के बीच कोई भी संख्या
पा सकते हैं। यदि अपने $(k+1)$वें पासे से $1$ फेंकें, तो वह कुल में $1$
जोड़ता है। अतः $k$ पासे फेंककर और फिर $(k+1)$वें पासे से $1$ फेंककर हम
$k+1$ और $6k+1$ के बीच कोई भी मान फेंक सकते हैं। क्या बचा? $6k+2$ से
$6k+6$ तक के मान। इन्हें हम पहले $k$ बार $6$ और फिर अपने $(k+1)$वें पासे
से $2$ और $6$ के बीच की संख्या फेंककर पा सकते हैं। मिलाकर इसका अर्थ है कि
$k+1$ पासों से $k+1$ और $6(k+1)$ के बीच कोई भी संख्या फेंक सकते हैं,
अर्थात् हमने मान्यता $P(k)$, आगमन परिकल्पना, का उपयोग करके $P(k+1)$ सिद्ध
कर दिया।
\end{proof}

बहुत बार हम आगमन का उपयोग तब करते हैं जब वस्तुओं (संख्याओं, समुच्चयों,
इत्यादि) की ऐसी शृंखला के बारे में कुछ सिद्ध करना हो जो स्वयं
``आगमनात्मक रूप से'' परिभाषित है, अर्थात् $(n+1)$वीं वस्तु को $n$वीं के
पदों में परिभाषित करके। उदाहरण के लिए $n$ तक की प्राकृतिक संख्याओं का योग
$s_n$ इस प्रकार परिभाषित कर सकते हैं
\begin{align*}
  s_0 & = 0\\
  s_{n+1} & = s_n + (n+1)
\end{align*}
यह परिभाषा देती है:
\begin{align*}
  s_0 & = 0,\\
  s_1 & = s_0 + 1 && = 1,\\
  s_2 & = s_1 + 2 && = 1 + 2 = 3\\
  s_3 & = s_2 + 3 && = 1 + 2 + 3 = 6, \text{ इत्यादि।}
\end{align*}
अब हम आगमन से सिद्ध कर सकते हैं कि $s_n=n(n+1)/2$।

\begin{prop}
  $s_n = n(n+1)/2$.
\end{prop}

\begin{proof}
  हमें सिद्ध करना है (1) $s_0=0\cdot(0+1)/2$ और (2) यदि
  $s_k=k(k+1)/2$, तो $s_{k+1}=(k+1)(k+2)/2$। (1) स्पष्ट है। (2) सिद्ध
  करने के लिए हम आगमन परिकल्पना $s_k=k(k+1)/2$ मानते हैं। इसका उपयोग करके
  हमें दिखाना है कि $s_{k+1}=(k+1)(k+2)/2$।

  $s_{k+1}$ क्या है? परिभाषा से $s_{k+1}=s_k+(k+1)$। आगमन परिकल्पना से
  $s_k=k(k+1)/2$। इसे पिछले समीकरण में प्रतिस्थापित कर सकते हैं, और फिर
  केवल भिन्नों का कुछ अंकगणित चाहिए:
  \begin{align*}
    s_{k+1} & = \frac{k(k+1)}{2} + (k+1) = {}\\
    & = \frac{k(k+1)}{2} + \frac{2(k+1)}{2} = {}\\
    & = \frac{k(k+1) + 2(k+1)}{2} = {}\\
    & = \frac{(k+2)(k+1)}{2}.
  \end{align*}
\end{proof}

यहाँ महत्त्वपूर्ण सीख यह है कि यदि आप किसी आगमनात्मक रूप से परिभाषित अनुक्रम
$a_n$ के बारे में कुछ सिद्ध कर रहे हैं, तो आगमन स्पष्ट विधि है। और यदि ऐसा
न भी हो (जैसे पासों के फेंक की संभावनाओं में), तब भी आप आगमन का उपयोग कर
सकते हैं यदि किसी प्रकार $k+1$ की स्थिति को $k$ की स्थिति से जोड़ सकें।

\end{document}
